\section{What does not follow}\label{sec:not}

The Companion grew for months before its statements were audited, and many of them were stated as theorems.  This section lists what
was withdrawn, corrected or found not to follow.  Readers who meet these statements in the Companion should read them with the status
given here; the theorem labels there do not describe the present status.  Each item names the layer that settled it.

\subsection{Status tags}

Every statement in this paper and in the audited Companion carries a status: \emph{proved} (with stated assumptions),
\emph{computed} (exact or numerical, with stated range), \emph{conditional} (holds under named choices), \emph{reading} (an
identification offered, not derived), \emph{open}, or \emph{withdrawn}.  Statements about physical quantities also name the frame in
which they are read: matter \([M]\), tachyon \([T]\) or light \([P]\).  Several of the errors below were comparisons across frames
that this tag would have exposed (Appendix~\ref{app:qam}, row~4).

\subsection{The Riemann source}

The Companion began from a separate paper that claimed a proof of the Riemann hypothesis and supplied the ``clock'' of the
framework: the zeros of \(\zeta\) as the points from which space-time was to be built.  That proof does not hold.  Since layer
v4.8.0-r04 the Riemann paper is no longer a source of proofs, and every statement that used it lost its justification:
\begin{center}
\small
\begin{tabular}{@{}>{\raggedright\arraybackslash}p{0.40\textwidth}>{\raggedright\arraybackslash}p{0.52\textwidth}@{}}
\toprule
statement & status\\
\midrule
unity of the three objects; unitary equivalence of the kernels & open (used Stone--von Neumann from the Riemann paper)\\
Bell consistency; strings and loop quantum gravity as compensator sectors & equivalence step open\\
Einstein equation from a trace formula; GUE statistics derived & open; the trace formula is not a defined action\\
Riemann hypothesis as a consequence & withdrawn; read as a conjecture only\\
Dirac operator on the antisymmetric sector (OP2) & existence and self-adjointness open\\
\bottomrule
\end{tabular}
\end{center}
The compensator, the cascade of Section~\ref{sec:octonions} and the triple hierarchy up to \(\OO\) were re-founded on Hilbert-space
projections and Cayley--Dickson multiplication and do not depend on it.  The open items are not refuted; they lost their proof.

\subsection{Statistics of the zeros}

That the spacing statistics of the Riemann zeros follow GUE is an empirical fact; the derivation of it from the framework is withdrawn
(layer v4.8.0-r05).  The claim that unfolding in the light frame singles out \(\mu=\gamma_L\) is not confirmed: the minimum lies at
\(\mu=1\), and the claim is contradicted on the first \(2000\) zeros (layer v4.8.0-r08).

\subsection{The kernel}

The Minkowski reading of the original norm was wrong (Remark~\ref{rem:signature}), and the sector rule by the sign of
\(\operatorname{Re}z\) was not Lorentz-invariant; both are replaced by Reading~B (layers v4.8.0-r06, r07).  The fibre \(S^7\) is not a group
containing \(G_2\), and the chain \(S^{15}\to S^8\to S^4\to S^2\) is not a composition of Hopf maps (Section~\ref{sec:octonions}).  The
statement that \(\operatorname{span}(e_1,e_2,e_3)\) carries the conjugate colour triplet confused two complex units
(Section~\ref{sec:generation}).  The spectral identification of the compensator inside the compensator matrix needs two hypotheses
that the Companion did not state (Proposition~\ref{prop:spectral}).

\subsection{The cosmological constant}

The original chain set \(\Lambda=0\) in the light frame and explained the observed value as an effect of the matter frame.  It does not
hold.
\begin{itemize}
\item The trace formula used as an action is not a defined action (layer v4.8.0-r05).
\item The value \(\beta_{\rm sym}^2/\ell_P^2\) exceeds the measured \(\Lambda\) by a factor of about \(10^{121}\), and no power of \(\gamma_L\)
closes the gap (layer v4.8.0-r03).
\item A corollary gave zero and a finite ratio \(\gamma_L^5\) for the same quantity; two remarks put \(\beta_{\rm sym}^2\) on \(\Lambda\) where a
theorem puts it on \(\Omega_M\) (layers r69, r67).
\item The regularization of the old balance was printed wrongly; the correct expansion is \(R(\tau)=-1/\tau+1-\tfrac{11}{24}\tau+\dots\) with
counterterm \(-\log(\varepsilon/\mu)\) and finite part \(1.686953\) at \(T=0.393\), \(\mu=1\) (layers r69, r70).
\end{itemize}
What survives is \(\Omega_M=\beta_{\rm sym}^2\) as a present-epoch reading (Section~\ref{sec:readouts}).

\subsection{Frame-shifted constants}

The Companion assigned each constant a ``frame-insertion count'' and derived a light-frame Planck length
\(\ell_P^{(P)}=\gamma_L^{-3/2}\ell_P^{(M)}\), a drift of \(\alpha\) as a frame effect, and a vacuum ratio \(\gamma_L^5\).  Under the counting
reading the frame difference lies in \(\gamma\), not in \(\hbar\) or \(G_N\), and these statements are not supported (layer v4.9.0-r64).  The
Hawking section assigned a tachyon counting to the outgoing photon, against its own frame convention (layers r64, r65).  Which of three
ways of converting constants between frames applies is left open by the author.

\subsection{The old predictions}

The original Companion listed six predictions P1--P6 derived from \(\beta_{\rm sym}\) and \(\gamma_L\).  None is confirmed (layers v4.8.0-r03,
r08, r09):
\begin{center}
\small
\begin{tabular}{@{}>{\raggedright\arraybackslash}p{0.34\textwidth}>{\raggedright\arraybackslash}p{0.58\textwidth}@{}}
\toprule
prediction & status\\
\midrule
P1 GUE improvement from light-frame unfolding & not confirmed; minimum at \(\mu=1\)\\
P2 Koide phase \(\theta_K=\pi-\sqrt2\arcsin\beta_{\rm sym}\) & excluded by the lepton masses, at \(300\)--\(400\sigma\)\\
P3 \(\beta_{\rm sym}=3^{-1/2}\exp(-\gamma_1/4\pi\gamma_L^5)\) & coincidence: exponent chosen after the fact, look-elsewhere about \(11\%\)\\
P4 \(\Lambda\) as a frame artefact & contradicted as stated (factor \(10^{121}\))\\
P5 annual ratio of \(\alpha\)- and \(G_N\)-drift & exponent corrected to \(\tfrac12\); not testable at present\\
P6 Hawking temperature in the tachyon frame & assumption made explicit; untested\\
\bottomrule
\end{tabular}
\end{center}

\subsection{Readings that went too far}

Two readings were corrected after independent recomputation (layer v4.9.0-r68).  A relation \(\gamma_M+\gamma_T=|\gamma_3|\) presented as a finding
holds by definition.  And the statement that the lepton data select the ``corrected'' horizon counting is not supported: they select the
U(1) counting, which the later literature considers less well founded (Section~\ref{sec:beta}).

\subsection{What is never claimed}

The Companion keeps a list of statements that are explicitly not made, and every layer checks against it: no exact finite closure of
the residual \(C_5\); no exact terminality of all readouts; no global closure \(C_6/T_2\); no global clock theorem and no Riemann theorem;
no global closure of all readouts.  The safe summary of the whole is: \emph{local, effective and frame-typed closures, not exact global
closure.}  This paper adds two items to that list.  The framework does not, at present, contain a dynamics: there is no action varied in
all fields from which its numbers follow.  And none of its numerical readouts is a derivation.

\medskip\noindent\emph{Status.}  This section is a register.  Each item is documented with its reasons in the layer named; the
inherited text of the Companion has not been rewritten, so the old statements remain there, marked with status notes.
